\subsection{A sharp cubic R\'enyi bound from binary grouping}
\label{sec:fqe-renyi-cubic-bound}

Binary grouping gives an independent entropy bound that will control the region of large moment slack. Set
\[
q_b=\frac{4\ln2}{4\ln2-1},\qquad
r_b=q_b-1=\frac1{4\ln2-1},\qquad
\frac32<q_b<\frac{313}{200}<\frac{1587}{1000}=q_s.
\tag{BT}\label{eq:fqe-binary-threshold}
\]
Thus $q_b/(r_b\ln2)=4$ and $1/2<r_b<1$.
The bounds on $q_b$ follow from three terms of the atanh series:
$56/81+2/1215<\ln2<1123/1620<3/4$.

\begin{lemma}[Binary grouping]
\label{lem:fqe-renyi-grouping}
For any finite probability vectors $u,v$ and $0\le a\le1$, set
$w=a^2(3-2a)$. Then
\[
H_{q_b}(au,(1-a)v)
\le4a(1-a)+wH_{q_b}(u)+(1-w)H_{q_b}(v).
\tag{BG}\label{eq:fqe-renyi-grouping}
\]
For $0<a<1$, equality holds exactly when $a=1/2$ and
$H_{q_b}(u)=H_{q_b}(v)$.
In particular $H_q(a,1-a)\le4a(1-a)$ holds for every $a$ if and only if
$q\ge q_b$; its equality cases are $a=0,1/2,1$.
\end{lemma}
\begin{proof}
Write $q=q_b$, $r=r_b$, $L=\ln2$ and
$h(w)=-w\ln w-(1-w)\ln(1-w)$.
The log-sum inequality gives, for positive $A,B$,
$\ln(A+B)\ge w\ln A+(1-w)\ln B+h(w)$.
Apply it with $A=a^q2^{-rH_q(u)}$ and
$B=(1-a)^q2^{-rH_q(v)}$. The required assertion follows if
\[
J(a):=q\{w\ln a+(1-w)\ln(1-a)\}+h(w)+4rL a(1-a)\ge0.
\]
The function $J$ is symmetric about $1/2$, and
$J(1/2)=J(1)=0$, with the latter understood by continuity.
Since $w-a=a(1-a)(2a-1)$ and $4rL=q$, differentiation cancels the
terms outside the derivative of $w$:
\[
J'(a)=6a(1-a)\left(q\ln\frac a{1-a}-\ln\frac w{1-w}\right).
\]
For $a>1/2$, put $t=a/(1-a)>1$. Then
\[
\frac w{1-w}=t\,T_2(t),\qquad T_2(t)=\frac{t(t+3)}{3t+1},
\qquad \operatorname{sgn}J'(a)=\operatorname{sgn}\{r\ln t-\ln T_2(t)\}.
\]
The function $s\mapsto\ln T_2(e^s)$ is strictly convex for $s>0$, since
\[
\left(t\frac d{dt}\right)^2\ln T_2(t)
=\frac{24t(t^2-1)}{(t+3)^2(3t+1)^2}>0.
\]
Its derivative starts at $1/2$ and tends to $1$.
Consequently $\ln T_2(e^s)-rs$ starts at zero with negative derivative,
is strictly convex, tends to $+\infty$, and has exactly one further zero.
Thus $J$ increases and then decreases between its zero endpoints $1/2,1$;
it is strictly positive in between. Symmetry proves $J\ge0$ throughout.
Equality in the log-sum inequality at $a=1/2$ requires $A=B$,
which is precisely the stated entropy condition.

Taking both conditional vectors pure proves the binary bound at $q_b$;
order monotonicity proves it for every $q\ge q_b$ with the same equality
cases. Necessity follows as $a=\epsilon\downarrow0$:
$H_q(\epsilon,1-\epsilon)=q\epsilon/((q-1)\ln2)+o(\epsilon)$,
whereas $4\epsilon(1-\epsilon)=4\epsilon+O(\epsilon^2)$.
Orders $q\le1$ are excluded by monotonicity from any $1<q<q_b$.
\end{proof}

\begin{proposition}[Sharp cubic entropy bound]
\label{prop:fqe-cubic-bound}
For every four-point probability vector $\lambda$ and every $q\ge q_b$,
\[
H_q(\lambda)\le\frac{10}3-6p+\frac83t
=4e_2(\lambda)+8e_3(\lambda),\qquad
p=\sum_i\lambda_i^2,\quad t=\sum_i\lambda_i^3.
\tag{CB}\label{eq:fqe-cubic-bound}
\]
The threshold $q_b$ is necessary for a bound valid on the entire simplex.
Equality occurs exactly at pure, uniform-two and uniform-four spectra.
\end{proposition}
\begin{proof}
By order monotonicity it suffices to take $q=q_b$.
Group the entries as
$\lambda=(au,a(1-u),(1-a)v,(1-a)(1-v))$.
For $0<a<1$, Lemma~\ref{lem:fqe-renyi-grouping} and its binary consequence give
\[
H_q(\lambda)\le4a(1-a)
+4a^2(3-2a)u(1-u)+4(1-a)^2(1+2a)v(1-v).
\]
The right side is exactly $4e_2+8e_3$, by expanding across the two groups.
If $a=0$ or $1$ the assertion is the binary bound.
The identity with $(p,t)$ follows from Newton's identities and $\sum\lambda_i=1$.
For equality with $0<a<1$, grouping forces $a=1/2$ and equal conditional
entropies; binary equality then makes both conditional spectra pure or
both uniform-two. This gives uniform-two or uniform-four. The boundary
case also permits a pure spectrum. Each such spectrum stays an equality
at every order, and all other spectra are strict already at $q_b$.
The binary restriction proves sharpness of the threshold.
\end{proof}

\begin{corollary}[Analytic cubic envelope]
\label{cor:fqe-renyi-high-f}
Every triple of four-point probability vectors satisfies, for every $q\ge q_b$,
\[
\mathcal E_q\le\frac53-\frac{20}{27}(P-1)
-\frac2{27}(F-1/2).
\tag{HC}\label{eq:fqe-renyi-high-f}
\]
In particular $\mathcal E_q\le5/3$ whenever $20(P-1)+2F\ge1$.
Equality in the $5/3$ ceiling holds exactly at a permutation of
$((1/2,1/2,0,0),(1/4)^4,(1/4)^4)$.
No quadratic variance bound or physical compatibility is required.
\end{corollary}
\begin{proof}
Sum \eqref{eq:fqe-cubic-bound} over the three cuts, substitute
$T=(17P-12-F)/12$, and divide by three.
Equality in the ceiling requires $20(P-1)+2F=1$ and equality in all three
scalar bounds. Each cut is then pure, uniform-two or uniform-four, with entropy
$0$, $1$ or $2$ at every order, and a mean of $5/3$ forces exactly one uniform-two and two
uniform-four spectra; these have $P=1$, $F=1/2$ and attain $5/3$ at every order.
\end{proof}
